% Part: first-order-logic
% Chapter: tableaux
% Section: provability-consistency

\documentclass[../../../include/open-logic-section]{subfiles}

\begin{document}

\iftag{FOL}
      {\olfileid{fol}{tab}{prv}}
      {\olfileid{pl}{tab}{prv}}

\olsection{\usetoken{S}{derivability} en consistentie}

We bewijzen nu enkele eigenschappen van de !!{derivability}srelatie.
Deze zijn op zichzelf interessant, maar spelen ook elk een rol in het
bewijs van de volledigheidsstelling.

\begin{prop}\ollabel{prop:provability-contr}
  Als $\Gamma \Proves !A$ en $\Gamma \cup \{!A\}$ inconsistent is,
  dan is $\Gamma$ inconsistent.
\end{prop}

\begin{proof}
Er zijn eindige verzamelingen $\Gamma_0 = \{!B_1, \dots, !B_n\}$ en
$\Gamma_1 =\{!C_1, \dots, !C_n\} \subseteq \Gamma$ zodat
  \begin{align*}
    \{\sFmla{\False}{!A}, &
    \sFmla{\True}{!B_1}, \dots, \sFmla{\True}{!B_n}\} \\
    \{\sFmla{\True}{!A}, &
    \sFmla{\True}{!C_1}, \dots, \sFmla{\True}{!C_m}\}
  \end{align*}
  gesloten !!{tableau}s hebben. Door de regel \Cut{} op $!A$ toe te
  passen, kunnen we deze combineren tot één gesloten !!{tableau} dat
  bewijst dat $\Gamma_0 \cup \Gamma_1$ inconsistent is. Omdat
  $\Gamma_0 \subseteq \Gamma$ en $\Gamma_1 \subseteq \Gamma$, geldt
  $\Gamma_0 \cup \Gamma_1 \subseteq \Gamma$; dus is $\Gamma$~inconsistent.
\end{proof}

\begin{prop}
\ollabel{prop:prov-incons}
$\Gamma \Proves !A$ dan en slechts dan als
$\Gamma \cup \{\lnot !A\}$ inconsistent is.
\end{prop}

\begin{proof}
Neem eerst aan dat $\Gamma \Proves !A$. Er is dan dus een gesloten
!!{tableau} voor
\[
\{\sFmla{\False}{!A},
\sFmla{\True}{!B_1}, \dots, \sFmla{\True}{!B_n}\}
\]
Door de regel $\TRule{\True}{\lnot}$ toe te passen, kunnen we dit
omzetten in een gesloten !!{tableau} voor
\[
\{\sFmla{\True}{\lnot !A},
\sFmla{\True}{!B_1}, \dots, \sFmla{\True}{!B_n}\}.
\]

Omgekeerd kunnen we een gesloten !!{tableau} voor de laatste
verzameling omzetten in een gesloten !!{tableau} voor de eerste.
Daarvoor verwijderen we elke formule die ontstaat door
\TRule{\True}{\lnot} toe te passen op de eerste aanname
$\sFmla{\True}{\lnot !A}$, evenals die aanname zelf, en voegen we de
aanname $\sFmla{\False}{!A}$ toe. Als een tak eerder gesloten was
omdat hij de conclusie bevatte van \TRule{\True}{\lnot} toegepast op
$\sFmla{\True}{\lnot !A}$, namelijk $\sFmla{\False}{!A}$, dan is de
overeenkomstige tak in het nieuwe !!{tableau} eveneens gesloten. Als
een tak in het oude tableau gesloten was omdat hij zowel de aanname
$\sFmla{\True}{\lnot !A}$ als $\sFmla{\False}{\lnot !A}$ bevatte,
kunnen we er een gesloten tak van maken door
$\TRule{\False}{\lnot}$ toe te passen op
$\sFmla{\False}{\lnot !A}$ en zo $\sFmla{\True}{!A}$ te verkrijgen.
Dit sluit de tak, omdat we $\sFmla{\False}{!A}$ als aanname hebben
toegevoegd.
\end{proof}

\begin{prob}
Bewijs dat $\Gamma \Proves \lnot !A$ dan en slechts dan als
$\Gamma \cup \{!A\}$ inconsistent is.
\end{prob}

\begin{prop}\ollabel{prop:explicit-inc}
  Als $\Gamma \Proves !A$ en $\lnot !A \in \Gamma$, dan is $\Gamma$
  inconsistent.
\end{prop}

\begin{proof}
  Stel dat $\Gamma \Proves !A$ en $\lnot !A \in \Gamma$. Dan zijn er
  $!B_1$, \dots, $!B_n \in \Gamma$ zodat \
  \[
  \{\sFmla{\False}{!A}, \sFmla{\True}{!B_1}, \dots, \sFmla{\True}{!B_n}\}
  \]
  een gesloten tableau heeft. Vervang de aanname
  \sFmla{\False}{!A} door \sFmla{\True}{\lnot !A} en voeg na de
  aannamen de conclusie toe van \TRule{\True}{\lnot}, toegepast op
  \sFmla{\False}{!A}. Elke !!{sentence} in het !!{tableau} die in het
  oude !!{tableau} met een verwijzing naar regel~$1$ werd
  gerechtvaardigd, wordt nu gerechtvaardigd met een verwijzing naar
  regel~$n+1$. Als het oude !!{tableau} gesloten was, is het nieuwe dat
  dus ook. Dit bewijst dat $\Gamma$ inconsistent is, want alle aannamen
  behoren tot~$\Gamma$.
\end{proof}

\begin{prop}\ollabel{prop:provability-exhaustive}
  Als $\Gamma \cup \{!A\}$ en $\Gamma \cup \{\lnot !A\}$ beide
  inconsistent zijn, dan is $\Gamma$ inconsistent.
\end{prop}

\begin{proof}
  Als er $!B_1$, \dots, $!B_n \in \Gamma$ en $!C_1$, \dots,
  $!C_m \in \Gamma$ zijn zodat
  \begin{align*}
    \{\sFmla{\True}{!A}, &
    \sFmla{\True}{!B_1}, \dots, \sFmla{\True}{!B_n}\} \text{ and}\\
    \{\sFmla{\True}{\lnot !A}, &
    \sFmla{\True}{!C_1}, \dots, \sFmla{\True}{!C_m}\}
  \end{align*}
  beide gesloten !!{tableau}s hebben, kunnen we één gecombineerd
  !!{tableau} construeren dat bewijst dat $\Gamma$ inconsistent is.
  We nemen daarvoor $\sFmla{\True}{!B_1}$, \dots,
  $\sFmla{\True}{!B_n}$ samen met $\sFmla{\True}{!C_1}$, \dots,
  $\sFmla{\True}{!C_m}$ als aannamen en passen vervolgens de regel
  \Cut{} toe. Hierdoor ontstaan twee takken: de ene begint met
  $\sFmla{\True}{!A}$ en de andere met $\sFmla{\False}{!A}$.  
  
  Voeg aan de linkerkant het deel van het eerste !!{tableau} onder de
  aannamen toe. Elke regeltoepassing blijft hier correct, omdat elke
  aanname van het eerste !!{tableau}, waaronder
  $\sFmla{\True}{!A}$, beschikbaar is. Elke tak onder
  $\sFmla{\True}{!A}$ sluit dus. 
  
  Voeg aan de rechterkant het deel van het tweede !!{tableau} onder
  zijn aanname toe, maar verwijder de resultaten van alle toepassingen
  van~$\TRule{\True}{\lnot}$ op $\sFmla{\True}{\lnot !A}$. De conclusie
  van $\TRule{\True}{\lnot}$ op $\sFmla{\True}{\lnot !A}$ is
  $\sFmla{\False}{!A}$. Deze is toch beschikbaar, want zij is de
  conclusie van de regel \Cut{} aan de rechterkant van het
  gecombineerde !!{tableau}.

  Als een tak in het tweede tableau gesloten was omdat hij zowel de
  aanname $\sFmla{\True}{\lnot !A}$ (die in het gecombineerde
  !!{tableau} niet langer als aanname voorkomt) als
  $\sFmla{\False}{\lnot !A}$ bevatte, kunnen we
  $\TRule{\False}{\lnot}$ toepassen op
  $\sFmla{\False}{\lnot !A}$ en zo $\sFmla{\True}{!A}$ verkrijgen. De
  overeenkomstige tak in het gecombineerde !!{tableau} sluit nu ook,
  want hij bevat de rechterconclusie van de regel \Cut{}, namelijk
  $\sFmla{\False}{!A}$. Als een tak in het tweede !!{tableau} om een
  andere reden sloot, sluit de overeenkomstige tak in het gecombineerde
  !!{tableau} eveneens. Elke !!{signed formula} behalve
  $\sFmla{\True}{\lnot !A}$ die op de tak in het oude, tweede
  !!{tableau} voorkomt, komt namelijk ook voor op de overeenkomstige tak
  in het gecombineerde !!{tableau}.
  \end{proof}

\end{document}
